% Presentation of original arrow notation. See flow-manifest.json.
\begin{center}
\begin{tikzpicture}[book flow node/.style={draw=BookHeading,line width=.65pt,fill=BookPanel,rounded corners=1.5pt,align=center,inner sep=5pt,minimum height=10mm,font=\fontsize{9}{11}\selectfont,execute at begin node={\hyphenpenalty=10000\exhyphenpenalty=10000}},book flow edge/.style={-{Latex[length=2mm]},draw=BookHeading,line width=.7pt}]
\node[book flow node,text width=39mm] (n0) at (0,0) {Rgba\allowbreak{}Frame};
\node[book flow node,text width=39mm] (n1) at (0,-1.8) {fit source pixels into target cell grid};
\node[book flow node,text width=39mm] (n2) at (0,-3.6) {sample brightness/\allowbreak{}alpha/\allowbreak{}color per cell or sub-cell};
\node[book flow node,text width=39mm] (n3) at (0,-5.4) {choose renderer};
\node[book flow node,text width=39mm] (n4) at (-5,-7.5) {charset: density ramp from custom characters};
\node[book flow node,text width=39mm] (n5) at (0,-7.5) {braille: 2x4 dot mask};
\node[book flow node,text width=39mm] (n6) at (5,-7.5) {quad: 2x2 quadrant mask};
\node[book flow node,text width=39mm] (n7) at (0,-9.8) {emit Bijou Surface};
\draw[book flow edge] (n0) -- (n1);
\draw[book flow edge] (n1) -- (n2);
\draw[book flow edge] (n2) -- (n3);
\draw[book flow edge] (n3) -- (n4);
\draw[book flow edge] (n3) -- (n5);
\draw[book flow edge] (n3) -- (n6);
\draw[book flow edge] (n4) -- (n7);
\draw[book flow edge] (n5) -- (n7);
\draw[book flow edge] (n6) -- (n7);
\end{tikzpicture}
\end{center}
